% ============================================================================
%  Chapter 21 -- Input and Output                                (Lecture 21)
% ============================================================================
\chapter{Input and Output}
\label{ch:io}

\epigraphline{Need input!}{Number 5, in the film \emph{Short Circuit}, John
Badham, 1986}

\begin{chapgoals}
By the end of this chapter you will be able to
\begin{itemize}
\item describe how a picture is formed on a CRT and on an LCD, and what a TFT active matrix adds;
\item explain the path of a key press: matrix, encoder, shift register, serial interface;
\item say what a mouse sends to the computer, and how often;
\item describe the stages of laser printing;
\item explain how the \prg kit reads its joystick, and compare polling, interrupts and DMA.
\end{itemize}
\end{chapgoals}
% Generated by tools/textbook/export_listings.js from tools/decks/mastery.js and tools/concepts/registry.json: do not edit.
\begin{chapstatus}{Learning status of this chapter}
\masterychip{now}{UNDERSTAND NOW}\enspace polling, interrupts, DMA; CRT and LCD displays; keyboard and mouse; printers.\par\smallskip
\masterychip{use}{USE NOW}\enspace the complete program, the PRG32 tools (used, not studied, in this course).\par\smallskip
\masterychip{later}{UNDERSTAND LATER}\enspace nothing.
\end{chapstatus}
\begin{chapstatus}{Roadmap of this chapter}
By the end of the chapter you can:
\begin{enumerate}[nosep,leftmargin=2em]
\item describe how displays, keyboards and printers work
\item compare polling with interrupts for reading a device
\end{enumerate}
It builds on: only the setup of Chapter~\ref{ch:setup}. Next: Chapter~\ref{ch:verm1p4}.
\end{chapstatus}


% ----------------------------------------------------------------------------
\section{Warm-up on the machine}

\asmfile[caption={Read the world, once per frame.},label=lst:io:warm]{code/l21/c03.S}

The reader has read the joystick since the first chapters
(Listing~\ref{lst:io:warm}). Each button is a switch on one input pin of the
chip: pressed or released, one bit. The runtime reads the pins, which are
memory-mapped registers configured as inputs, and packs them into one word: a
decoded representation, in the sense of Chapter~\ref{ch:combinational}. And
the game asks once per frame: this is called \term{polling}.

The listing also repeats a lesson of Chapter~\ref{ch:instructions}. The mask
of the buttons is kept in \reg{s0}, a saved register, and therefore survives
the call that plays the note: the marker moves with \key{LEFT} also while
\key{A} sounds. In \reg{t0} the mask would be lost at the call. (The program
was built as a cartridge and run on the headless runner of \prgqt.)

A keyboard has a hundred keys, a mouse moves in two dimensions, a display has
thousands of pixels: how do they talk to a processor?

% ----------------------------------------------------------------------------
\section{Displays}

\subsection{The video: an introduction}

The first displays showed mostly characters, in black and white; today they
are graphical and in colour. A graphical display is a matrix of points, the
\term{pixels}, each lit with a different intensity, and an image is composed
by switching the pixels on and off. A different density of pixels gives a
different resolution; the colours go from 256 (8 bits) to 16 million (24
bits) and beyond. The size is measured in inches along the diagonal. The
display of the \prg kit has a diagonal of 2.8 inches, 320 by 240 pixels and
65\,536 colours.

\subsection{The cathode-ray tube}

\begin{table}[htb]
\centering\small
\caption{A short history of the cathode-ray tube.}\label{tab:io:crt}
\begin{tabularx}{\textwidth}{@{}llL@{}}
\toprule
\textbf{Year} & \textbf{Who} & \textbf{What}\\
\midrule
early 1870s & William Crookes & the Crookes tube: the first cathode-ray tube, a laboratory instrument\\
1897 & Karl Ferdinand Braun & the Braun tube and the first cathode-ray oscilloscope: the first practical application\\
1907 & Boris Rosing & a CRT displays simple geometric shapes from a received video signal: the root of television\\
1934 & Telefunken & the first television set with a CRT on the market, in Germany\\
1970s--1990s & the computer industry & the CRT monitor: the display of every terminal and personal computer\\
\bottomrule
\end{tabularx}
\end{table}

In a \term{cathode-ray tube} (CRT; Table~\ref{tab:io:crt},
Figure~\ref{fig:io:crt}) the cathode is a small metal element heated to
incandescence, which emits electrons by the thermionic effect. Inside the
tube the electrons are directed into a beam by a high difference of potential
between cathode and anode, and other electromagnetic fields focus the beam.
The beam, the ``electron brush'', is deflected by magnetic fields (the
Lorentz force) or by electric fields, and strikes a point on the inner
surface of the screen, which is coated with phosphor and lights up. Sweeping
the beam line by line, many times per second, paints the picture.

\begin{figure}[htb]
\centering
\begin{subfigure}[b]{0.3\textwidth}\centering\includegraphics[width=\textwidth,height=36mm,keepaspectratio]{mac_classic.jpg}\caption{A compact Macintosh: a 9-inch CRT}\end{subfigure}\hfill
\begin{subfigure}[b]{0.36\textwidth}\centering
\begin{tikzpicture}[x=0.42cm,y=0.42cm,font=\sffamily\tiny]
  \draw[prgsteel,thick,fill=prgpanel] (0,-0.9) -- (4.2,-0.9) -- (8,-2.9) -- (8.5,-2.9) arc (-35:35:5.05) -- (8,2.9) -- (4.2,0.9) -- (0,0.9) -- cycle;
  \draw[prgred,thick] (0.5,-0.25) rectangle (1.3,0.25);
  \fill[prgsteel] (2.0,1.1) rectangle (3.4,1.6) (2.0,-1.6) rectangle (3.4,-1.1);
  \fill[prgamber] (4.4,1.5) rectangle (5.6,2.0) (4.4,-2.0) rectangle (5.6,-1.5);
  \draw[prggreen,very thick,dashed] (1.3,0) -- (4.6,0) .. controls (6,0.2) .. (9.35,1.6);
  \fill[prggreen] (9.35,1.6) circle (2.5pt);
  \node[anchor=north] at (0.9,-1.0) {cathode};
  \node[anchor=south] at (2.7,1.6) {focusing};
  \node[anchor=south] at (5.0,2.0) {deflection};
  \node[anchor=west] at (9.6,1.6) {phosphor};
  \node[anchor=north] at (5.2,-3.0) {electron beam, anode at high voltage};
\end{tikzpicture}
\caption{The scheme of a CRT}\end{subfigure}\hfill
\begin{subfigure}[b]{0.3\textwidth}\centering\includegraphics[width=\textwidth,height=36mm,keepaspectratio]{crt_scope.jpg}\caption{An oscilloscope used as a display}\end{subfigure}
\caption{The cathode-ray tube.}
\label{fig:io:crt}
\end{figure}

\subsection{Liquid crystal displays}

LCD monitors have almost completely replaced the cathode-ray tube: they give
a better picture, with less bulk, weight and power, and less strain on the
eyes. \term{Liquid crystals} are viscous organic molecules that flow like a
liquid but have a spatial structure. An electric field changes the
orientation of the molecules, and with it an optical property of the crystal:
the angle of polarisation of the light that crosses it. A crystal is a valve
for light, opened and closed by a voltage.

\subsection{Inside an LCD panel}

The screen is made of many cells, often one per pixel, in which the liquid
crystals are trapped (Figure~\ref{fig:io:lcd}). Each cell has electrical
contacts, so that a field can be applied to the liquid that it contains. The
cells lie between two polarising filters whose axes are perpendicular to each
other, with a light, the \term{backlight}, behind the rear one. Polarisation
decides whether the light passes: when the crystals are not driven the light
is rotated and crosses the front polariser; when they are driven it is
completely blocked. In a colour screen each cell is divided into three
sections with a red, a green and a blue filter: the three fields of the
RGB565 colour of Chapter~\ref{ch:information} end up here, as three voltages.

\begin{figure}[htb]
\centering
\begin{subfigure}[b]{0.46\textwidth}\centering\includegraphics[width=\textwidth,height=38mm,keepaspectratio]{lcd_layers.jpg}\caption{The layers of an LCD}\end{subfigure}\hfill
\begin{subfigure}[b]{0.46\textwidth}\centering\includegraphics[width=\textwidth,height=38mm,keepaspectratio]{lcd_pixels.jpg}\caption{Red, green and blue sub-pixels}\end{subfigure}
\caption{Liquid crystal displays.}
\label{fig:io:lcd}
\end{figure}

\subsection{Driving the cells: from segments to the active matrix}

In calculators and digital watches there is one electrical contact per
segment, which becomes impossible as the number of elements grows. A
\emph{passive matrix} has one group of contacts per row and one per column,
not per pixel; only one pixel at a time is driven, and the others must
remember their state. The consequences are \emph{persistence}, the trails
behind moving objects, and crosstalk between neighbouring cells.

High-resolution screens use an \term{active matrix}. The panel contains a
thin film of transistors (TFT, \emph{thin-film transistor}) on a transparent
semiconductor layer; the field is applied exactly where it is needed by a
transistor in each cell, and the transistor stores the electrical state of
its pixel while the others are updated. Each pixel is, in effect, an analogue
memory cell. The display of the kit is a TFT panel with an ILI9341
controller. In the more recent OLED displays each pixel emits its own light,
and no backlight is needed.

% ----------------------------------------------------------------------------
\section{Keyboard and mouse}

\subsection{The keyboard}

The keyboard derives from that of typewriters, and is the main device for the
input of characters. Its keys are grouped: alphanumeric keys; special keys
such as enter and backspace; arrow keys and function keys. Keyboard and video
are not directly connected to each other: it is the processor that reproduces
on the video the information typed on the keyboard. The echo of a key on the
screen is a program: read a code, draw a character. Today's keyboards
connect through USB, or without wires.

\subsection{The path of a key press}

\begin{figure}[htb]
\centering
\begin{tikzpicture}
  \foreach \k/\x/\y in {Q/0/0,W/1/0,E/2/0,R/3/0,A/0.3/1,S/1.3/1,D/2.3/1,F/3.3/1,Z/0.6/2,X/1.6/2,C/2.6/2,V/3.6/2}{
    \def\kd{D}\ifx\k\kd \def\c{prgamber}\else\def\c{prgsteel}\fi
    \node[draw=white,fill=\c,text=white,rounded corners=1pt,minimum size=5.5mm,font=\ttfamily\bfseries\small] at (\x*0.62,-\y*0.62) {\k};
  }
  \node[lbl,text=prgcomment] at (1.2,-1.9) {\emph{rows and columns of switches}};
  \node[blkacc,text width=27mm,minimum height=14mm] (e) at (5.2,-0.6) {\textbf{encoder}\\\footnotesize row and column to an $n$-bit code};
  \node[blkvio,text width=27mm,minimum height=14mm] (s) at (8.6,-0.6) {\textbf{shift register}\\\footnotesize $n$ bits in parallel to $n$ in series};
  \node[blkgrn,text width=27mm,minimum height=14mm] (u) at (12.0,-0.6) {\textbf{serial interface}\\\footnotesize e.g.\ a USB controller};
  \draw[->,thick] (2.9,-0.6) -- (e); \draw[->,thick] (e) -- (s); \draw[->,thick] (s) -- (u);
\end{tikzpicture}
\caption{The path of a key press.}
\label{fig:io:key}
\end{figure}

The keys of a keyboard are switches arranged in rows and columns
(Figure~\ref{fig:io:key}). A key that is pressed joins one row to one column;
an \emph{encoder} turns that combination of row and column into a key code of
$n$ bits; a \emph{shift register} turns the $n$ bits from parallel to serial;
and a \emph{serial interface}, for example a USB controller, sends them to
the computer.

\begin{prgconcept}
Every block of Figure~\ref{fig:io:key} is a machine that the reader has
already designed. This is the pay-off of Chapters \ref{ch:combinational}
to~\ref{ch:registers}.
\end{prgconcept}

\subsection{The mouse}

The mouse was invented in the 1960s by Douglas Engelbart (the first prototype
is of 1964, the patent was filed in 1967 and the famous public demonstration
is of 1968): it was the first pointing device in history. It was a wooden box
with two wheels underneath and a wire leaving from the back, like the tail of
a mouse. The two wheels detected the movements along the horizontal and the
vertical axis, and the electronics turned them into electrical signals. In a
\emph{mechanical} mouse a ball turns two rollers, each of which drives a
counter; in an \emph{optical} mouse a source of light and a tiny camera
measure the movement of the surface underneath.

\subsection{What a mouse sends}

Every time it has moved by a minimum distance, about a hundredth of an inch,
called a \emph{mickey}, the mouse sends three bytes: how many units it has
moved along $x$ since the last report; the movement along $y$; and the state
of the buttons. The two movements are signed numbers in two's complement
(Chapter~\ref{ch:integers}). Modern mice report up to a thousand times per
second. Inside, the mouse is two counters, a shift register that puts the
three bytes in series, and a serial interface; on the side of the computer
the control software, the \term{device driver}, supervises the communication
with the processor.

% ----------------------------------------------------------------------------
\section{Printers}

\subsection{The printer}

A printer is an output device that prints information on paper. Its speed is
measured in lines per minute, characters per second or pages per minute; its
resolution, that is, the quality of the print, in dots per inch. The
daisy-wheel and chain printers are now disused; a \emph{dot-matrix} printer
strikes an inked ribbon with needles; an \emph{ink-jet} printer throws tiny
drops of ink; and then there is the laser printer. A printer is also a small
computer: it has its own processor, memory and firmware.

\subsection{The laser printer}

The laser printer derives from the photocopier
(Figure~\ref{fig:io:laser}). It works in six stages.

\begin{enumerate}
\item A drum is charged electrically, to about 1\,000 volts.
\item A laser beam, reflected by a rotating mirror, scans the length of the
  drum; the beam is modulated so as to produce a pattern of light and dark
  points.
\item The points struck by the beam lose their charge. After each line the
  drum turns by a fraction of a degree.
\item The toner is attracted by the points that are still charged: an image of
  the line is formed.
\item The drum is pressed against the paper and transfers the toner to it;
  heated rollers fix it permanently.
\item The drum is discharged and cleaned.
\end{enumerate}

The page is first built as a bitmap in the memory of the printer: a frame
buffer for paper.

\begin{figure}[htb]
\centering
\includegraphics[width=0.5\textwidth]{laser_printer.jpg}
\caption{The laser writes the page on the drum, line by line.}
\label{fig:io:laser}
\end{figure}

% ----------------------------------------------------------------------------
\section{Back to the machine}

\subsection{The devices of the PRG32 kit}

\begin{table}[htb]
\centering\small
\caption{The devices of the \prg kit.}\label{tab:io:kit}
\begin{tabularx}{\textwidth}{@{}P{0.2\textwidth}lLL@{}}
\toprule
\textbf{Device} & \textbf{Direction} & \textbf{How it is connected} & \textbf{How a cartridge uses it}\\
\midrule
joystick and buttons & input & one switch per input pin (GPIO) & \code{prg32\_input\_read}: one bit per button\\
display, $320\times 240$ TFT & output & SPI serial bus, with DMA & the \code{prg32\_gfx\_*} calls draw into the frame buffer\\
amplifier and speaker & output & I2S serial bus & the \code{prg32\_audio\_*} calls: notes, samples, music\\
Wi-Fi radio & input and output & inside the chip & cartridge upload, the Store, multiplayer\\
USB serial port & input and output & USB & flashing the firmware, debug messages\\
\bottomrule
\end{tabularx}
\end{table}

One detail of Table~\ref{tab:io:kit} was promised in
Chapter~\ref{ch:electronics}: a switch bounces for a few milliseconds when it
is pressed, and sampling it once per frame, every few tens of milliseconds,
hides the bounce.

\subsection{Three ways to deal with a device}

\begin{description}[style=nextline,leftmargin=1.5em]
\item[Polling (programmed I/O).] The program asks the device, again and again,
  whether it has something. It is simple, and it is what a game does once
  per frame; it is wasteful when nothing happens.
\item[Interrupts.] The device asks for attention when it needs it; the
  processor suspends what it is doing, serves the device and resumes
  (Chapter~\ref{ch:instructions}).
\item[Direct memory access.] The device exchanges whole blocks of data with
  the memory by itself; the processor only starts the transfer and is told
  when it ends (Chapter~\ref{ch:memory}).
\end{description}

\begin{prgmachine}
On the kit all three are at work at the same time: the game polls the
joystick, the timer that paces the frames interrupts, and the display and the
sound are fed by DMA.
\end{prgmachine}

% ----------------------------------------------------------------------------
% Generated by tools/textbook/export_listings.js from decks/plan.js and code/l21.S: do not edit.
\section{The program of this chapter}\label{sec:lab:l21}

The listings of this chapter are parts of one program, \fname{code/l21.S} in the
repository of the course, printed in full in Listing~\ref{lst:lab:l21}. The
automatic tests of the course material build it and run it on the headless
runner of \prgqt{}, with scripted input, and check the frame;
Figure~\ref{fig:lab:l21} is the frame that the test wrote.

\begin{description}[style=nextline,leftmargin=1.5em]
\item[On the screen.] A yellow square in the middle of the screen. LEFT moves it to the left; A plays a note, and the square keeps moving while the note sounds.
\item[What to press.] LEFT moves the square; A plays a note.
\end{description}

\begin{figure}[htb]
\centering
\includegraphics[width=0.5\textwidth]{lab_l21.png}
\caption{\fname{l21.S} on \prgqt{}.}\label{fig:lab:l21}
\end{figure}

\asmfile[caption={The complete program of the chapter, \fname{code/l21.S}.},label=lst:lab:l21]{code/l21/program.S}

To run it on \prgqt{} the commands are those of Listing~\ref{lst:lab:l21:run};
Section~\ref{sec:prg:testing} gives the same steps for QEMU and for the board.

\begin{shell}[caption={Building \fname{l21.S} and running it on \prgqt{}.},label=lst:lab:l21:run]
# 1. start PRG32-QT on this computer (it listens on port 8080)

# 2. in the PRG32 folder, with ESP-IDF active: build the cartridge
python3 -m prg32 cartridge build ~/arc12/code/l21.S --portable \
    --entry-prefix l21 --name l21 --out l21.prg32

# 3. send it to PRG32-QT and start it
python3 -m prg32 esp32c6 upload-and-run l21.prg32 \
    --url http://localhost:8080 --slot cart0

# 4. look, press, compare with the previous slide
\end{shell}


% ----------------------------------------------------------------------------
% Generated by tools/textbook/export_listings.js from tools/concepts/assessments.json: do not edit.
\begin{chapstatus}{The thread of this chapter}
Reading the joystick once per frame is polling: simple, and blind between two reads. Next: the finished game.
\par\medskip\textbf{Exit ticket} (two to four minutes, alone, on paper):
\begin{enumerate}[nosep,leftmargin=2em]
\item In four steps, how does a laser printer turn a page into toner on paper? Why does a printer have a processor and memory of its own?
\item A game reads the joystick once per frame, at 30 frames per second. A tap lasts 20 ms and falls between two reads. Is it seen? What would an interrupt change?
\end{enumerate}
\end{chapstatus}


\section{Check your understanding}

\begin{quiz}{What does the encoder of a keyboard controller do?}
\item It lights the key that was pressed
\item It turns a row and column combination into a key code
\item It sends the code to the monitor
\item It counts how long a key is held
\end{quiz}

\begin{quiz}{In a colour LCD, a pixel is made of:}
\item one cell that changes colour
\item three cells with red, green and blue filters
\item a phosphor dot hit by electrons
\item a transistor that emits light
\end{quiz}

\begin{quiz}{Why does a passive-matrix LCD show trails behind moving objects?}
\item Its backlight is too weak
\item Its pixels are driven one at a time and must hold their state in between
\item It has too few colours
\item The video card is slow
\end{quiz}

\begin{quiz}{A \prg game reads the joystick once per frame. This technique is called:}
\item direct memory access
\item interrupt-driven input
\item polling
\item memory mapping
\end{quiz}

\section{Exercises}

\begin{exercise}[paper first, then \prg]
\begin{enumerate}[nosep]
\item A keyboard has 104 keys: how many bits does the key code need? How many
  rows and columns could the matrix have?
\item A page of 8 by 11 inches at 600 dots per inch, one bit per dot: how much
  memory does the printer need for the bitmap?
\item Move a pointer on the \prg screen with the joystick, as a mouse would:
  keep $x$ and $y$, and clamp them to the screen.
\item Make button \key{A} act only once per press (edge detection,
  Chapter~\ref{ch:sequential}).
\item Draw a seven-segment score with rectangles, using the transcoder table
  of Chapter~\ref{ch:combinational}.
\end{enumerate}
\expected{On paper: 7 bits and, for example, 8 rows by 13 columns;
$4\,800\times 6\,600$ dots, that is 31\,680\,000 bits, about 4~MB. On \prg: a
pointer that stays on the screen, and a button that fires once per press.}
\end{exercise}

\section{Glossary}

\begin{glossarytable}
cathode-ray tube & tubo a raggi catodici\\
liquid crystals & cristalli liquidi\\
backlight & retroilluminazione\\
active matrix & matrice attiva\\
keyboard / key & tastiera / tasto\\
pointing device & periferica di puntamento\\
dot-matrix / ink-jet printer & stampante ad aghi / a getto d'inchiostro\\
drum / toner & tamburo / toner\\
polling & interrogazione ciclica\\
device driver & driver di periferica\\
\end{glossarytable}

\section{Conclusions and reading}

Input/output devices are a fundamental component of computers. Display
technology has made enormous progress, from the electron beam to a transistor
behind every pixel. Keyboard, mouse and printer have changed little in form,
and inside they are built from encoders, counters and shift registers. A
device reaches the processor through a controller and a bus, and is handled
by polling, by interrupts or by DMA. On the \prg kit the whole chain, from a
switch to a pixel, can be followed instruction by instruction.

\paragraph{Reading.} \citet{tanenbaum2013}, section~2.4 (input/output);
\citet{stallings2021}, the chapter on input/output; \emph{Playing with
RISC-V}, chapter~6, ``State and Input'' \citep{prg32textbook}.

The fourth block is complete. The next chapter finishes \vermone, and the one
after it publishes the game.
